\section{Proof dependencies and the physical source}
\label{app:v95-dependencies}
The original problem and the three input statements preceding the new
proofs fix every source used in this article. The collision source is
$\one_\Pi\nu_R^*$ and its angular density is
$\omega/(4\pi c_*\chi)$. No independent boundary spectral theorem is
assumed.

The two-ended packet theorem uses the Borel segment disintegration
(module 196), the full angular-density derivative (modules 180--185),
the original guard derivative (177), the one-sided trace theorem (220),
and the positive-cutoff volume estimate audited in 196 and 222.
Its endpoint minimum is new; the historical inlet-volume and
outlet-volume comparisons remain separately credited to 198 and 223.

The intrinsic anchor uses the same disintegration and radial growth,
and the critical-volume estimate from 187 in the explicit form
\eqref{eq:v94-critical-finite}. Its shell layer cake is restated and
proved in \eqref{eq:v95-shell-input}; that calculation is inherited
from 224. The target-adapted two-sided anchor and relatively thin
classification are the new use of these inputs.

The bidirectional comparison uses the fixed-count signed-prefix
identity of 221 and the projected Stieltjes calculation of 225.
The forward calculation is written in full and includes its positive
terminal and prefix cuts. Standard BV terminology follows \cite{AFP}.
The scalar conclusion uses only the original fixed-width source bound
and positive scalar remainder stated in Input~\ref{input:v94-scalar}.
No path-remainder theorem enters the new upper estimate.

The main article directly loads only modules 226--228. The complete
archive retains modules 1--225 and the full preceding article, with
historical reference namespaces. Preservation is not an independent
reconstruction of the inherited four-record spectral or coarea proofs.
